\subsection{The complete quartic character calculation for two old coordinates}

\begin{statement}{A.96}{Proved}
{Quartic common fields and two-coordinate initialization}
For nonzero rational numbers $D_1,\ldots,D_k$, put
\[
 \mathcal E_4=\Q[Z_1^{\pm1},\ldots,Z_k^{\pm1}]
                         /(Z_i^4-D_i),\qquad
 K=\Q(\sqrt{D_1},\ldots,\sqrt{D_k}).
 \tag{A.96.1}
\]
For every nonempty subset $J\subset\{1,\ldots,k\}$ such that
$\prod_{i\in J}D_i$ is a negative rational square, let $r_J>0$
be the rational number determined by
$\prod_{i\in J}D_i=-r_J^2$.
The largest abelian field contained in every field factor
of $\mathcal E_4$ is exactly
\[
              B_4=K\bigl(\sqrt{2r_J}:
                       \prod_{i\in J}D_i=-r_J^2\bigr).
 \tag{A.96.2}
\]
The empty list of additional generators gives $B_4=K$.
In particular every common character is quadratic.

Now let $a,b\in\Q^\times$, let
\[
        C:\ aX^4-bY^4=1,\qquad T=aX^4,\qquad
        c\in\Q\setminus\{0,1\},\qquad
        D_0=c/a,\quad D_1=(c-1)/b.
 \tag{A.96.3}
\]
For every common character $\chi$ of the fibre algebra
$\Q[Z_0^{\pm1},Z_1^{\pm1}]/(Z_0^4-D_0,Z_1^4-D_1)$,
the cyclic class $(T-c,\chi)$ on $C$ is the restriction of
an explicitly computed constant in $\operatorname{Br}(\Q)[2]$.
Its evaluation at every ordinary local point with $T\ne c$
is that same constant's local evaluation.

Consequently the finite-character test of A.85 always vanishes
for two old finite punctures at exponent four.
More explicitly, for an everywhere locally soluble old quartic
cover with two finite punctures, any new rational puncture, any
finite set $S$ of places, and one prescribed ordinary old local point
at each $v\in S$ avoiding that puncture admit a rational new coefficient whose augmented
quartic cover is everywhere locally soluble and whose
forgetful map lifts every prescribed point.
This assertion concerns the actual finite initialization
problem at exponent four.
\end{statement}
\begin{proof}
\emph{The finite Galois maps and all their common characters.}
Write
\[
 D_i=(-1)^{\delta_i}\prod_{P\in\mathcal P}P^{w_{iP}},
       \qquad \delta_i\in\{0,1\},\quad w_{iP}\in\Z.
\]
We use the
actual fields and restriction map of A.57:
\[
 \begin{gathered}
 F=\Q(\mu_8,\sqrt P:P\in\mathcal P),\qquad
 L=F(P^{1/4}:P\in\mathcal P),\\
 \pi:G=\operatorname{Gal}(L/\Q)\twoheadrightarrow
                     \Gamma=\operatorname{Gal}(F/\Q).
 \end{gathered}
 \tag{A.96.4}
\]
The field $F$ is multiquadratic.
The proof of (A.57.7) gives $[G,G]=\operatorname{Gal}(L/F)$.
Thus every abelian character trivial on every root-field
stabilizer factors through this specific $\Gamma$.
If $\Omega$ is the set of geometric root tuples and
$H=\langle\pi(G_\omega):\omega\in\Omega\rangle$, its common
abelian field is $F^H$.
The square roots $Z_i^2$ give $K\subset F^H$ and
$H\subset\Gamma_K=\operatorname{Gal}(F/K)$.
If $i\notin K$, A.80 already proves $H=\Gamma_K$.
There is then no negative-square subset relation, because
such a relation would put a square root of $-1$ in $K$.
This proves (A.96.2) in that case.

Suppose henceforth that $i\in K$.
For $\sigma\in\Gamma_K$ write
\[
 \sigma(\zeta_8)=\zeta_8^u,\quad u\in\{1,5\},\qquad
 \sigma(\sqrt P)=(-1)^{\epsilon_P}\sqrt P,
                         \quad\epsilon_P\in\{0,1\}.
 \tag{A.96.5}
\]
The restriction $u\in\{1,5\}$ follows because $\sigma$ fixes
$i$.
All lifts of $\sigma$ to $G$ have
$P^{1/4}\mapsto\zeta_4^{k_P}P^{1/4}$ with
$k_P=\epsilon_P+2h_P$ and arbitrary $h_P\in\mathbf F_2$,
by the independent Kummer degree calculation in A.57.
This retains the actual relations between $u$ and the
$\epsilon_P$ in $F$; no independent cyclotomic signs
have been added.
The root tuples are
$z_i=\zeta_8^{\delta_i+2e_i}\prod_P P^{w_{iP}/4}$.
Their literal stabilizer equations are
\[
 (u-1)(\delta_i+2e_i)+2\sum_Pw_{iP}k_P=0\pmod8.
 \tag{A.96.6}
\]
Since $4\mid u-1$, the variables $e_i$ disappear from
these equations.
Also $\sum_Pw_{iP}\epsilon_P$ is even: this is precisely
the condition that $\sigma$ fix $\sqrt{D_i}$ once it
fixes $i$.
Consequently (A.96.6) is equivalent to the following
ordinary linear system over $\mathbf F_2$:
\[
 \sum_P\bar w_{iP}h_P
   =-\frac12\sum_Pw_{iP}\epsilon_P
        -\frac{u-1}{4}\delta_i\pmod2
                       \qquad(1\le i\le k).
 \tag{A.96.7}
\]
Here $\bar W=(\bar w_{iP})$ is the reduction modulo two
of the integral valuation matrix $W$.

Let $t=(t_i)\in\mathbf F_2^k$ be a relation
$t^t\bar W=0$, and put
\[
 \delta_t=\sum_i t_i\delta_i\pmod2,\qquad
 r_t=\prod_P P^{\frac12\sum_i t_iw_{iP}}>0,\qquad
 \beta_t=2^{\delta_t}r_t.
 \tag{A.96.8}
\]
The exponents defining $r_t$ are integers, so $r_t\in\Q^\times$.
The corresponding subset product of the $D_i$ is
$(-1)^{\delta_t}r_t^2$.
The linear system (A.96.7) has a solution exactly when its
right side pairs trivially with every such $t$.
That pairing is
\[
       \sum_P\left(\frac12\sum_i t_iw_{iP}\right)\epsilon_P
                  +\frac{u-1}{4}\delta_t
                        =\chi_{\beta_t}(\sigma)\pmod2.
 \tag{A.96.9}
\]
The first summand is the character of $\sqrt{r_t}$.
For $u=1,5$, the second is $\delta_t$ times the character
of $\sqrt2$, since
$\sigma(\sqrt2)=(-1)^{(u-1)/4}\sqrt2$.
Thus (A.96.9) is an equality of actual quadratic characters
on $\Gamma_K$.

It follows that the union of the images of all root
stabilizers is already the subgroup
\[
 \mathcal A
     =\Gamma_K\cap\bigcap_{t^t\bar W=0}\ker\chi_{\beta_t}.
 \tag{A.96.10}
\]
Indeed, solving (A.96.7) supplies one actual lift
$(\sigma,(k_P)_P)$ fixing a root tuple through (A.96.6);
conversely every fixed tuple satisfies those equations.
Consequently $H=\mathcal A$, rather than merely a
subgroup generated by a possibly smaller uncomputed set.
Galois correspondence gives
$F^H=K(\sqrt{\beta_t}:t^t\bar W=0)$.

We now eliminate the positive-square relations from this
last expression.
Since $i\in K$, there is a relation $t_0$ with
$\delta_{t_0}=1$.
For any relation $t$ with $\delta_t=0$, the vector
$t'=t+t_0$ is a relation with $\delta_{t'}=1$.
If $J$ is the overlap of the subsets represented by
$t$ and $t_0$, the exact positive rational identity is
\[
             r_t r_{t_0}
                    =r_{t'}\prod_{i\in J}|D_i|.
 \tag{A.96.11}
\]
Modulo rational squares, it expresses $r_t$ as
$(2r_{t_0})(2r_{t'})\prod_{i\in J}|D_i|$.
The square roots of the last factors belong to $K$,
because $K$ contains every $\sqrt{D_i}$ and also $i$.
Hence all the positive-relation generators already
belong to the field obtained from the negative ones.
This proves (A.96.2).
There is also a direct map witnessing each negative
generator inside every field factor: if
$z=\prod_{i\in J}z_i$, then $z^4=-r_J^2$,
$\iota=z^2/r_J$ has square $-1$, and
\[
                         ((1-\iota)z)^2=2r_J.
 \tag{A.96.12}
\]
This is compatible with the common-field calculation,
including every possible root choice.

\emph{The individual characters and the negative singleton relations.}
On (A.96.3), the two mod-two Steinberg identities of A.80 are
\[
             (T-c,D_0)=(a,D_0),\qquad
             (T-c,D_1)=(b,D_1).
 \tag{A.96.13}
\]
For example $(T-c)/a+D_0=X^4$; division by this square,
the Steinberg relation and bilinearity give the first
identity, and the other uses $Y$ and $b$.
If $D_0=-r^2$, then
\[
 \begin{split}
 T-c&=a(X^4+r^2),\\
 X^4+r^2
      &=(X^2+r)^2-2rX^2\\
      &=N_{\Q(\sqrt{2r})/\Q}(X^2+r+\sqrt{2r}\,X).
 \end{split}
 \tag{A.96.14}
\]
The norm is interpreted in the quadratic etale algebra
when $2r$ is a square.
It gives $(T-c,2r)=(a,2r)$.
For the other negative singleton one obtains in exactly
the same variables
$T-c=b(Y^4+r^2)$ and $(T-c,2r)=(b,2r)$.
Thus both possible singleton generators have explicit
constant classes.

\emph{The negative two-coordinate relation.}
It remains to treat
$D_0D_1=-r^2$, equivalently $c(c-1)=-ab r^2$.
Define the nonzero rational numbers
\[
                    u=a(1-c),\qquad k=\frac{br}{1-c}.
 \tag{A.96.15}
\]
The old curve equation and this exact relation give
\[
 \begin{split}
 T-c&=u(X^4+k^2Y^4),\\
 X^4+k^2Y^4
       &=(X^2+kY^2)^2-2k(XY)^2\\
       &=N_{\Q(\sqrt{2k})/\Q}
                    (X^2+kY^2+\sqrt{2k}\,XY),\\
 \frac{2k}{2r}&=\frac{b}{1-c}=\frac{D_0}{r^2}.
 \end{split}
 \tag{A.96.16}
\]
For the first line, its coefficients are
$a(1-c)$ and $bc$, so the right side is
$(1-c)T+c(T-1)=T-c$.
The final line records the exact squareclass comparison
between the two quadratic fields used in the calculation.
Taking quaternion symbols, and then applying (A.96.13),
gives the explicit constant
\[
              (T-c,2r)
                 =(u,2k)+(a,D_0)
                      \quad\text{in }\operatorname{Br}\Q(C).
 \tag{A.96.17}
\]
All common characters in (A.96.2), for two radicands,
are sums of the characters handled in
(A.96.13)--(A.96.17).
Adding the corresponding constants proves the Brauer
assertion for every character.
The identities are defined and specialize at every
old local point with $X,Y,T-c\ne0$.
In particular their assertion is about the actual local
evaluations, not only about vanishing residues.

\emph{The exact affine comparison and the initialization map.}
For two arbitrary distinct old finite punctures
$c_1,c_2$, write their old equations as
$T-c_1=a_1X^4$, $T-c_2=a_2Y^4$.
The change of variables
\[
 d=c_2-c_1,\qquad
 t=\frac{T-c_1}{d},\qquad
 a=\frac{a_1}{d},\quad b=\frac{a_2}{d},\quad
 c'=\frac{c-c_1}{d}
 \tag{A.96.18}
\]
is an isomorphism of the old covers, with inverse
$T=c_1+dt$.
It gives $aX^4-bY^4=1$, preserves the two fibre
radicands, and induces the literal cyclic-class identity
\[
                       (T-c,\chi)=(d,\chi)+(t-c',\chi).
 \tag{A.96.19}
\]
Thus the constant-class result holds for the original
punctures as well.

Choose the finite set $S$ in A.85, enlarged to contain
the prescribed places and the ramification of the
finitely many explicit constant generators above.
At any newly included place choose an old ordinary point
avoiding $c$; such a point exists by local solubility and
smoothness, since the excluded fibre is finite.
For each common character $\chi$, its evaluations at
the chosen points are the restrictions of one global
constant class.
Outside $S$ these evaluations vanish by A.85, or directly
by the chosen ramification condition.
Global reciprocity therefore gives
\[
       \sum_{v\in S\cup\{\infty\}}
            \operatorname{inv}_v(T(P_v)-c,\chi)=0.
 \tag{A.96.20}
\]
The exact sufficiency direction of A.85 now produces
$a_{\rm new}\in\Q^\times$ and the augmented cover
$T-c=a_{\rm new}Z^4$.
Its forgetful morphism
$(T,X,Y,Z)\mapsto(T,X,Y)$ lifts each prescribed local
point, and the augmented cover has points at every
place.
This proves the claimed initialization result.
\end{proof}

\paragraph{Scope of the calculation and retained exact searches.}
Equation (A.96.2) computes every common character for
arbitrarily many quartic radicands.
The constancy proof uses the fact that with two
radicands a negative-square subset has at most two
members.
For three members the nonconstant example in A.69 is
retained, and A.90 gives its completed realization.
At larger exponents the exact character calculation
is A.57; no higher-exponent constancy follows merely
from (A.96.2).
The proof also supplies no single rational new
coefficient preserving an arbitrarily prescribed
infinite collection of old local points.

The certificate
\path{code/quartic_common_character_certificate.py}, with
retained output
\path{code/quartic_common_character_certificate_output.txt},
checks the actual finite stabilizer equations, the generated
common-field quotient, its negative-relation generators, and
the polynomial and squareclass identities above.
The earlier exact searches
\path{code/tripod_brauer_initial_search.py} and
\path{code/tripod_joint_square_search.py}, with their retained
output files, motivated (A.96.14)--(A.96.17).
Their finite absence of witnesses is not used to prove
the general identities.
